\def\GraphicsFolder{appendices/superexchange-virtual-hopping/pictures}

\chapter{Superexchange and virtual hopping in Hubbard lattices}\label{app:superexchange-virtual-hopping}
 
Superexchange is a key mechanism in the formation of AF ordering in Hubbard lattices. It is an extension of the direct exchange effect, that we describe in Sec.~\ref{appsec:direct-exchange}, that favour antiferromagnetic ordering and takes advantage of the chemical structure of cuprates. The mechanism becomes clear enough if we start with a $2$-sites Hubbard toy model.

\section{Direct exchange and virtual hopping in Hubbard lattices}\label{appsec:hubbard-2}

Consider the $2$ sites Hubbard model of Eq.~\eqref{eq:hubbard-2-sites},
\[
	\hat H = 
	-t \left(
		 \hat c_{1\uparrow}^\dagger \hat c_{2\uparrow} + \hat c_{1\downarrow}^\dagger \hat c_{2\downarrow} + \hc 
	\right)
	+ U \left(
		 \hat n_{1\uparrow} \hat n_{1\downarrow} + \hat n_{2\uparrow} \hat n_{2\downarrow} 
	\right),
\]
with $i=1,2$ the site index. For an half-filled system (i.e. with two electrons) the Hilbert space is six-dimensional. We use the notation $\ket{n_{1\uparrow} n_{1\downarrow} n_{2\uparrow} n_{2\downarrow}}$ to indicate the six states of the computational basis:
\begin{equation}\label{appeq:hubbard-2-states}
	\begin{aligned}
		\ket{\psi_1} &\equiv \ket{1010} \\
		\ket{\psi_2} &\equiv \ket{1100} \\
	\end{aligned}
	\qquad
	\begin{aligned}
		\ket{\psi_3} &\equiv \ket{1001} \\
		\ket{\psi_4} &\equiv \ket{0110} \\
	\end{aligned}
	\qquad
	\begin{aligned}
		\ket{\psi_5} &\equiv \ket{0011} \\
		\ket{\psi_6} &\equiv \ket{0101} \\
	\end{aligned}
\end{equation}
We give a pictorial representation of these states in Fig.~\ref{appfig:two-sites-hubbard-model}. We ordered the states in such a way that the two states with total spin different from zero, $\ket{\psi_1}$ and $\ket{\psi_6}$, are at the extrema of the set.

\begin{figure}
	\centering
	\begin{tikzpicture}[
		site/.style={circle, fill=#1, minimum size=0.3em, outer sep=0pt, inner sep=0pt},
		up/.style={site=tabblue},
		down/.style={site=tabred},
		empty/.style={site=lightgray},
		double/.style={site=black},
		arrow/.style={-stealth, line width=1pt, shorten <=2pt, shorten >=2pt},
	]	
	
	% up -- up (psi1)
	\draw[lightgray] (0,0) grid (1,0);
	
	\node[up] (1-1up) at (0,0) {};
	\draw[tabblue, arrow] ([shift={(0,-0.3)}]1-1up.center) -- ([shift={(0,0.3)}]1-1up.center);
	
	\node[up] (1-2up) at (1,0) {};
	\draw[tabblue, arrow] ([shift={(0,-0.3)}]1-2up.center) -- ([shift={(0,0.3)}]1-2up.center);
	
	\node (psi1) at (2,0) {$\ket{\psi_1}$};
	
	% up down -- (psi2)
	\draw[lightgray] (0,-1) grid (1,-1);
		
	\node[double] (2-1double) at (0,-1) {};
	\draw[tabblue, arrow] ([shift={(-0.1,-0.3)}]2-1double.center) -- ([shift={(-0.1,0.3)}]2-1double.center);
	
	\draw[tabred, arrow] ([shift={(0.1,0.3)}]2-1double.center) -- ([shift={(0.1,-0.3)}]2-1double.center);
	
	\node[empty] (2-1double) at (1,-1) {};
	
	\node (psi2) at (2,-1) {$\ket{\psi_2}$};
	
	% up -- down (psi3)
	\draw[lightgray] (0,-2) grid (1,-2);
	
	\node[up] (3-1up) at (0,-2) {};
	\draw[tabblue, arrow] ([shift={(0,-0.3)}]3-1up.center) -- ([shift={(0,0.3)}]3-1up.center);

	\node[down] (3-2down) at (1,-2) {};
	\draw[tabred, arrow] ([shift={(0,0.3)}]3-2down.center) -- ([shift={(0,-0.3)}]3-2down.center);
		
	\node (psi3) at (2,-2) {$\ket{\psi_3}$};
	
	%% SECOND COLUMN %%
	
	% down -- up (psi4)
	\draw[lightgray] (4,0) grid (5,0);
	
	\node[down] (4-1down) at (4,0) {};
	\draw[tabred, arrow] ([shift={(0,0.3)}]4-1down.center) -- ([shift={(0,-0.3)}]4-1down.center);
	
	\node[up] (4-2up) at (5,0) {};
	\draw[tabblue, arrow] ([shift={(0,-0.3)}]4-2up.center) -- ([shift={(0,0.3)}]4-2up.center);
	
	\node (psi4) at (6,0) {$\ket{\psi_4}$};
	
	% -- up down (psi5)
	\draw[lightgray] (4,-1) grid (5,-1);

	\node[empty] (5-1empty) at (4,-1) {};

	\node[double] (5-2double) at (5,-1) {};
	\draw[tabblue, arrow] ([shift={(-0.1,-0.3)}]5-2double.center) -- ([shift={(-0.1,0.3)}]5-2double.center);
	
	\draw[tabred, arrow] ([shift={(0.1,0.3)}]5-2double.center) -- ([shift={(0.1,-0.3)}]5-2double.center);
		
	\node (psi5) at (6,-1) {$\ket{\psi_5}$};
	
	% down -- down (psi6)
	\draw[lightgray] (4,-2) grid (5,-2);
	
	\node[down] (6-1down) at (4,-2) {};	
	\draw[tabred, arrow] ([shift={(0,0.3)}]6-1down.center) -- ([shift={(0,-0.3)}]6-1down.center);
	
	\node[down] (6-2down) at (5,-2) {};
	\draw[tabred, arrow] ([shift={(0,0.3)}]6-2down.center) -- ([shift={(0,-0.3)}]6-2down.center);
	
	\node (psi3) at (6,-2) {$\ket{\psi_6}$};
\end{tikzpicture}
	\caption{Computational basis of Eq.~\eqref{appeq:hubbard-2-states} for the $2$-sites Hubbard model.}
	\label{appfig:two-sites-hubbard-model}
\end{figure}

The matrix elements of the hamiltonian can be computed using this basis,
\begin{equation}\label{appeq:hubbard-2-Hij}
	H_{ij} = \mel{\psi_i}{\hat H}{\psi_j}
	\quad\implies\quad
	H = \begin{bmatrix}
		0 &  0 &  0 &  0 &  0 &  0 \\
		0 &  U & -t & -t &  0 &  0 \\
		0 & -t &  0 &  0 & -t &  0 \\
		0 & -t &  0 &  0 & -t &  0 \\
		0 &  0 & -t & -t &  U &  0 \\
		0 &  0 &  0 &  0 &  0 &  0 \\  
	\end{bmatrix}
\end{equation}
In the above matrix, rows and columns are ordered following the states indices of Eq.~\eqref{appeq:hubbard-2-states}. The states $\ket{\psi_1}$ (both spins up) and $\ket{\psi_6}$ (both spins down) are zero-energy eigenstates. Indeed, they do not interact via Coulomb repulsion (are on different sites) and neither through NNs hopping, due to Pauli exclusion principle.

All states with null total spin projection, namely $\lbrace \ket{\psi_i} \rbrace$ for $2 \le i \le 5$ have null matrix elements with the aforementioned states, but non-zero matrix elements among themselves. Note that this set includes both the states with one doubly occupied site and one empty site ($\ket{\psi_2}$, $\ket{\psi_5}$) and those with two singly occupied sites ($\ket{\psi_3}$, $\ket{\psi_4}$). These states are connected by the central $4 \times 4$ matrix in Eq.~\eqref{appeq:hubbard-2-Hij}. The latter is diagonalised by the means of a change of basis $M$,
\[
	M \begin{bmatrix}
		 U & -t & -t &     \\
		-t &    &    & -t  \\
		-t &    &    & -t  \\
		   & -t & -t &  U  \\
	\end{bmatrix} M^\dagger = \begin{bmatrix}
	E^- &   & 	& 	 \\
	    & 0 & 	&	 \\
		&   & U & 	 \\
		&	&	& E^+
	\end{bmatrix}
\]
where blank slots indicate zeros. In
Tab.~\ref{apptab:two-sites-hubbard-model-eigenstates} we report the eigenvectors and relative eigenvalues obtained from diagonalisation. 

\begin{table}
	\centering
	\begin{tabular}{c c c}
		\textbf{Structure} & \textbf{Eigenstate} & \textbf{Energy} \\
		\midrule
		Spin-$1/2$ singlet & $\displaystyle \frac{\ket{\psi_3} - \ket{\psi_4}}{\sqrt{2}}$ & $\displaystyle E^- = \frac{U}{2} - \sqrt{\frac{U^2}{4} + 4t^2}$ \\
		&&\\ % Extra space
		Spin-$1/2$ triplet & $\ket{\psi_1}, \displaystyle \frac{\ket{\psi_3} + \ket{\psi_4}}{\sqrt{2}}, \ket{\psi_6}$ & 0 \\
		&&\\ % Extra space
		&$\displaystyle \frac{\ket{\psi_2}-\ket{\psi_5}}{\sqrt{2}}$ & $U$\\
		&&\\ % Extra space
		&$\displaystyle \frac{\ket{\psi_2}+\ket{\psi_5}}{\sqrt{2}}$ & $\displaystyle E^+ = \frac{U}{2} + \sqrt{\frac{U^2}{4} + 4t^2}$ \\
		\midrule
	\end{tabular}
	\caption{List of exact eigenstates and relative energies for the $2$-sites half-filled Hubbard model.}
	\label{apptab:two-sites-hubbard-model-eigenstates}
\end{table}

The ground-state is the singlet state,
\[
	\frac{\ket{\psi_3} - \ket{\psi_4}}{\sqrt{2}} = \frac{\ket{1010} - \ket{0101}}{\sqrt{2}} = \frac{\ket{\uparrow\downarrow} - \ket{\downarrow\uparrow}}{\sqrt{2}}
\]
of energy
\[
	E^- = \frac{U}{2} - \sqrt{\frac{U^2}{4} + 4t^2} \simeq - \frac{4t^2}{U}
\]
the latter approximation being true if $U \gg t$ (strong coupling limit). The singlet state pairs with a spatially-symmetric (node-less) wavefunction, while triplet states with an antisymmetric wavefunction (which, instead, has a node). The entire triplet (second row of Tab.~\ref{apptab:two-sites-hubbard-model-eigenstates}) has zero energy. Excited states are anti-symmetrised and symmetrised version of the polarised states $\ket{\psi_1}$ and $\ket{\psi_6}$.

\subsection{Direct exchange in Hubbard lattices}\label{appsec:direct-exchange}

The fact that we limited to two sites allowed for an exact diagonalisation of the toy hamiltonian of Eq.~\eqref{eq:hubbard-2-sites}. When moving to lattices, the above procedure cannot be extended easily. By using perturbation theory, instead, we can obtain the same result and easily extend to larger lattices.

Consider Eq.~\eqref{eq:hubbard-2-sites} separating two parts,
\[
	\hat H = \hat H_0 + \hat H_1
\]
where
\[	
\begin{aligned}
	\hat H_0 &\equiv U \left(
	\hat n_{1\uparrow} \hat n_{1\downarrow} + \hat n_{2\uparrow} \hat n_{2\downarrow} 
	\right) \\
	\hat H_1 &\equiv -t \left(
	\hat c_{1\uparrow}^\dagger \hat c_{2\uparrow} + \hat c_{1\downarrow}^\dagger \hat c_{2\downarrow} + \hc 
	\right)
\end{aligned}
\]
At zeroth order in perturbation theory, thus considering just the Coulomb repulsion $\hat H_0$, the Hilbert space separates in two manifolds
\[
\begin{aligned}
	\mathrm{Span}(\ket{\psi_1},\ket{\psi_3},\ket{\psi_4},\ket{\psi_6}) \qquad & E = 0 \\
	\mathrm{Span}(\ket{\psi_2},\ket{\psi_5}) \qquad & E = U
\end{aligned}
\]
We limit to the ground manifold, thus ``excluding from the Hilbert space'' the states with double occupation $\ket{\psi_2}, \ket{\psi_5}$.

First order corrections to energy are null:
\[
	\forall i,j
	\quad\colon\quad
	\mel{\psi_i}{\hat H_1}{\psi_j} = 0
\]
because the hopping operator $\hat H_1$ produces a doubly occupied site when applied to $\ket{\psi_3}, \ket{\psi_4}$ and zero when applied to $\ket{\psi_1}, \ket{\psi_6}$ due to Pauli exclusion principle.

We need to move to second order corrections. Let us change basis, and define the singlet state
\[
	\ket{s} \equiv \frac{\ket{\psi_3} - \ket{\psi_4}}{\sqrt{2}} 
\]
and the triplet states
\[
	\ket{t_+} \equiv \ket{\psi_1}
	\qquad
	\ket{t_0} \equiv \frac{\ket{\psi_3} + \ket{\psi_4}}{\sqrt{2}} 
	\qquad
	\ket{t_-} \equiv \ket{\psi_6}
\]
Let also $E_\psi^{(0)}$ be the zeroth-order energy associated to state $\ket{\psi}$. Second order correction to the singlet state energy is given by
\[
\begin{aligned}
	E_s^{(2)} &= \sum_{\psi \neq s} \frac{\mel{s}{\hat H_1}{\psi}\mel{\psi}{\hat H_1}{s}}{E_s^{(0)} - E_\psi^{(0)}} \\
	&= \frac{2t^2}{E_s^{(0)} - E_{\psi_2}^{(0)}} + \frac{2t^2}{E_s^{(0)} - E_{\psi_5}^{(0)}} = - \frac{4t^2}{U}
\end{aligned}
\]
The above result can be obtained easily by direct computation. Note that the $+$ sign in the middle is obtained as a combination of the internal $-$ sign in the definition of the singlet state, and another $-$ sign arising from Fermi anticommutation rules. 

Instead, for the central triplet state
\[
\begin{aligned}
	E_{t_0}^{(2)} &= \sum_{\psi \neq t_0} \frac{\mel{t_0}{\hat H_1}{\psi}\mel{\psi}{\hat H_1}{t_0}}{E_{t_0}^{(0)} - E_\psi^{(0)}} \\
	&= \frac{2t^2}{E_{t_0}^{(0)} - E_{\psi_2}^{(0)}} - \frac{2t^2}{E_{t_0}^{(0)} - E_{\psi_5}^{(0)}} = 0
\end{aligned}
\]
This time, the central $-$ sign is present precisely because of the internal sign of the state $\ket{t_0}$. Finally,
\[
	E_{t_\pm}^{(2)} = \sum_{\psi \neq t_0} \frac{\mel{t_\pm}{\hat H_1}{\psi}\mel{\psi}{\hat H_1}{t_\pm}}{E_{t_\pm}^{(0)} - E_\psi^{(0)}} = 0
\]
again, because of Pauli exclusion principle.

\begin{figure}
	\centering
	\begin{tikzpicture}[
	site/.style={circle, fill=#1, minimum size=0.3em, outer sep=0pt, inner sep=0pt},
	up/.style={site=tabblue},
	down/.style={site=tabred},
	empty/.style={site=lightgray},
	double/.style={site=black},
	arrow/.style={-stealth, line width=1pt, shorten <=2pt, shorten >=2pt},
	]	
	
	% up -- down (psi3)
	\draw[lightgray] (0,0) grid (1,0);
	
	\node[up] (3-1up) at (0,0) {};
	\draw[tabblue, arrow] ([shift={(0,-0.3)}]3-1up.center) -- ([shift={(0,0.3)}]3-1up.center);
	
	\node[down] (3-2down) at (1,0) {};
	\draw[tabred, arrow] ([shift={(0,0.3)}]3-2down.center) -- ([shift={(0,-0.3)}]3-2down.center);
	
	%% SECOND COLUMN %%
	
	% -- up down (psi5)
	\draw[lightgray] (2,0) grid (3,0);
	
	\node[empty] (5-1empty) at (2,0) {};
	
	\node[double] (5-2double) at (3,0) {};
	\draw[tabblue, arrow] ([shift={(-0.1,-0.3)}]5-2double.center) -- ([shift={(-0.1,0.3)}]5-2double.center);
	
	\draw[tabred, arrow] ([shift={(0.1,0.3)}]5-2double.center) -- ([shift={(0.1,-0.3)}]5-2double.center);
	
	%% THIRD COLUMN %%
	
	% up -- down (psi3)
	\draw[lightgray] (4,0) grid (5,0);
	
	\node[up] (3-1up) at (4,0) {};
	\draw[tabblue, arrow] ([shift={(0,-0.3)}]3-1up.center) -- ([shift={(0,0.3)}]3-1up.center);
	
	\node[down] (3-2down) at (5,0) {};
	\draw[tabred, arrow] ([shift={(0,0.3)}]3-2down.center) -- ([shift={(0,-0.3)}]3-2down.center);
	
\end{tikzpicture}
	\caption{Schematisation of the second order virtual hoppings for singlet and triplet states. \todo}
	\label{appfig:hubbard-2-virtual}
\end{figure}

What makes $\ket{s}$ and $\ket{t_0}$ fundamentally different? The singlet state takes advantage of the second-order virtual process where one spin hops to the neighbouring site and returns back. This process is schematised in Fig.~\ref{appfig:hubbard-2-virtual}, and is referred to as ``virtual hopping''. But each one of the two electrons can do this. The antisymmetry properties of many-body electronic wavefunctions imply, in the end, that these two virtual processes interfere constructively. For the triplet state $\ket{t_0}$, the two process are perfectly plausible but interfere destructively. \todo

\subsection{Superexchange in cuprates}

The scheme of direct exchange explains the insurgence of antiferromagnetism in the Hubbard model at strong coupling, $U \gg t$. Direct exchange \todo